% =====================================================================
% Prefazione: come si legge questo documento
% =====================================================================
% copre: docs/31-glossario.md
% verificato-al-commit: f41fd4c
% ---------------------------------------------------------------------

Questo lavoro nasce da una cartuccia rotta e da una curiosità che non passa.

La cartuccia è una copia originale di Pokemon Smeraldo, giocata da bambino e conservata per vent'anni, la cui tasca degli strumenti a un certo punto ha smesso di avere senso: da uno slot in avanti gli oggetti ordinari sono diventati Poke Ball. Non è un problema che si risolve rigiocando, perché quel salvataggio contiene un pezzo di tempo che non si rifà. Da qui la prima domanda, che è concreta e non teorica: si può aprire un salvataggio di vent'anni fa, capire che cosa è andato storto e ripararlo senza distruggere tutto il resto.

La curiosità è più vecchia della cartuccia rotta, ed è la seconda origine di questo documento. I giochi Pokemon delle prime tre generazioni hanno una proprietà che da ragazzo si intuisce e da ingegnere si riconosce: sono sistemi che comunicano fra loro attraverso un filo, e quel filo trasporta creature. Un Pokemon che passa da una console a un'altra è, sotto la superficie, una struttura di ottanta byte che attraversa un canale seriale sincrono, viene cifrata, verificata con un checksum e ricomposta all'altro capo. Da bambino si vede la creatura; dopo cinque anni di telecomunicazioni si vede il canale. Questo documento nasce dal desiderio di vedere entrambe le cose insieme, e di capire fino in fondo la macchina che c'è dietro un gesto che da piccolo sembrava magia.

Da queste due origini sono nati cinque lavori paralleli, che sono i casi di studio di questo testo. Riparare l'inventario di quella cartuccia agendo sul salvataggio estratto fisicamente. Costruire un ponte che porti i Pokemon dalle prime due generazioni alla terza su hardware originale, cosa che Nintendo non ha mai reso possibile e che quindi va costruita. Modificare un Nintendo 3DS per estrarre i dati delle proprie cartucce. Far parlare un computer con una Nintendo Switch attraverso il suo protocollo wireless proprietario, per scambiare Pokemon con un gioco che gira sulla console. E studiare come si automatizza un gioco pilotando una console e leggendone lo schermo.

Sono cinque obiettivi diversi, e nessuno di essi era all'inizio un esercizio accademico: ciascuno è nato da una cosa che volevo fare e che non si poteva fare senza capire come funzionava. È questo che tiene insieme il documento, e che spiega perché ogni capitolo teorico finisca in un vincolo pratico e ogni vincolo pratico risalga a un principio.

Detto questo, il testo che segue non è un diario di quei cinque lavori: è la spiegazione, dal primo bit, di che cosa si è dovuto imparare per farli. L'oggetto della trattazione è come si rappresenta, si verifica, si trasmette e si converte un dato strutturato fra sistemi digitali che non sono stati progettati per parlarsi. I Pokemon sono il carico utile, le console sono i nodi, e la scelta di questo dominio non è nostalgia travestita: quei giochi sono sistemi reali, completamente documentati da decompilazioni fedeli, abbastanza piccoli da poter essere capiti fino in fondo e abbastanza vecchi da non avere alcuna protezione che offuschi i meccanismi. Sono, in altre parole, un laboratorio migliore di quasi qualunque sistema moderno, e il fatto che siano anche i giochi della mia infanzia è la ragione per cui ho avuto la pazienza di leggerne i disassemblati.

\section*{Da quale zero si parte}

Da zero. La prima parte non dà per scontato che il lettore sappia che cos'è un bit, in
che senso un numero possa essere negativo dentro un byte, che cosa distingua un
indirizzo da un valore o che cosa significhi che un processore esegue qualcosa. Chi sa
tutto questo può saltarla, ma non è stata scritta per completezza formale: il resto
del documento vi si appoggia esplicitamente, e in più punti la spiegazione di un
fenomeno complicato si riduce a ricordare un fatto elementare enunciato là.

Dove serve un concetto più avanzato, lo si introduce nel punto in cui serve e non
prima. Quando quel concetto ha un nome consolidato nella teoria delle
telecomunicazioni, quel nome viene detto, perché riconoscere che un problema
apparentemente idiosincratico è un caso particolare di qualcosa di noto è metà del
lavoro. Alcuni esempi di ciò che il lettore troverà inquadrato in questo modo: il
checksum come codice di rilevazione d'errore, la conversione delle statistiche fra due
generazioni come quantizzazione con saturazione, lo stallo fra due protocolli come
problema di controllo di flusso, e il collegamento fra due console come canale sincrono
con un solo clock e nessun controllo d'errore.

\section*{Ogni affermazione ha una fonte, e le fonti hanno un ordine}

Ogni affermazione tecnica di questo documento porta una citazione, e le voci citate sono
le stesse registrate nel registro delle fonti del progetto. Non è una formalità
accademica: è la difesa contro un errore specifico che questo progetto ha commesso e
documentato. Quattro affermazioni su cui si stava per costruire del codice si sono
rivelate sbagliate, e sono emerse soltanto leggendo i sorgenti dei giochi invece delle
enciclopedie.

Ne è seguita una gerarchia, che vale in tutto il documento. Al primo livello stanno le
decompilazioni fedeli dei giochi e la documentazione dell'hardware: sono la definizione
del comportamento, e dove divergono da qualunque altra fonte hanno ragione loro. Al
secondo le wiki di dominio, utili per orientarsi e da verificare prima di appoggiarvi
una riga di codice. Al terzo le implementazioni di riferimento, che dimostrano che una
cosa è possibile ma non che sia il modo corretto di farla. Al quarto articoli, blog e
ricerca applicata. Al quinto forum e community, dove l'informazione è spesso l'unica
esistente e va trattata come testimonianza.

Quando un'affermazione non è verificata, non viene scritta come se lo fosse: sta dentro
un riquadro che la dichiara tale. Il lettore che trova un riquadro sa che quel punto è
un debito, non un risultato.

\section*{Che rapporto ha con il resto del progetto}

Questo documento non è l'unico testo del progetto, e la divisione dei compiti è
deliberata. Le note sotto \file{docs/} sono il percorso di studio: trattano gli stessi
argomenti per un lettore che ha già il contesto, e sono navigabili come grafo. La
referenza dei formati è la documentazione byte per byte, e serve a chi deve scrivere
codice adesso. Le schede sotto \file{.claude/} sono stato del lavoro, non conoscenza, e
non hanno un corrispondente qui.

Tre testi che dicono la stessa cosa in tre registri diversi divergono, se nulla lo
impedisce. Ciò che lo impedisce, in questo progetto, è un controllo che fallisce: ogni
capitolo dichiara in testa quali documenti copre e a quale commit è stato verificato
contro di essi, e lo strumento \file{tools/check-thesis-coverage.py} segnala i capitoli
il cui documento di riferimento è cambiato dopo l'ultima verifica, i documenti che
nessun capitolo copre, e le citazioni che non corrispondono ad alcuna voce di
bibliografia. La bibliografia stessa non è scritta a mano: è generata dalla stessa
tabella che genera le note sulle fonti, così che le due non possano dire cose diverse.

\section*{Convenzioni}

I valori esadecimali si scrivono nella forma \hx{2A}, i binari nella forma \bin{00101010}.
I percorsi dei file del progetto sono in \file{questo stile}. Un termine tecnico alla sua
prima occorrenza è in \term{corsivo}, e i termini raccolti nel glossario del progetto
mantengono la stessa definizione qui. Gli estratti di codice riportano nella didascalia
il file e la fonte da cui vengono, perché un estratto senza provenienza non è
verificabile.
